% Pursuing Stacks (cotes 134-2 to 134-8): archive cote 134-6, site page 34 (PDF page 35), typescript page 352.
% typescript. Draft compared: ps5.tex:3017-3072.
% AI-transcribed by Claude Opus 5.5 (claude-opus-5-5, adaptive thinking, effort high), Message Batches API, batch ps3 (msgbatch_01XdFmuVVBAcgc7EFxdtCqsj), 2026-11-12; not yet checked by eye.
\scanpage{134-6/34}
\mspage{352}

true for an arbitrary topos:  for instance the classifying topos $B_G$
of a discrete group $G$ has the second property, but\inserted{visibly} not the first
\deleted{visibly} unless $G$ is the unit group. (I recall that the category of
sheaves on $B_G$ is the category  (G-sets)  of sets on which $G$ operates.)
I feel both properties for a topos merit a name.\xout{xx}  The first \xout{xxxxxx}
(every sheaf has a section)  can be
viewed as the strongest conceivable (I would think) \inserted[typed]{global} asphericity proper-%
ty for a topos,  as far as $H^1$ goes at any rate,  as the $H^1$ of $X$  wi th
coefficients in \uline{any} group object will be zero.  (But I confess I did'
\margin{$<$} nt try and look if any precontractor,  say,  is aspheric ...) The
second property (every subtopos is \deleted{one of the} trivial \deleted{two}) comes with
a \xout{xxxx} rather different flavor, it suggest the image of just one
``point''  -  and as a matter of fact,  the etale topos of a scheme, say,
has this property iff it is reduced to a point.  Such a topos may be
called ``punctual'' (not to be confused though with some other meanings
suggested by this word,  such as being equivalent to the topos defined
by a one-point topological space,  namely \hand{$\mathcal{A}$} being equivalent to (Sets)\hand{)}
or ``atomic'' (which has rather unpleasant connotations though nowadays!\hand{),}
or maybe ``vertical'' (this image is suggested by the $B_G$ above) -  the
``base'' i.e. the final object of \hand{$\mathcal{A}$}  being very ``small'' (in terms of
harbouring subobjects),  so the inner structure is expressed like a
kind of tower,  related (in the case of $B_G$)  to the ``Galois tower''
of subgroups of $G$ ...  The corresponding notion of a ``horizontal''
topos is visibly the one when \hand{$\mathcal{A}$}  admits the subobjects of the final
object as a generating family.  In terms of these definitions,  A topos
is horizontal and vertical, iff it is either the``empty''or the ``final''
\hand{(}or ``one point''\hand{)}topos. This brings to mind that in the notion of verti-%
cality,  we should exclude the``empty''topos (which formally satisfies
the condition -  every sheaf has a section).  This brings to my atten-%
tion toothat I certainly do not want to consider \xout{the}\inserted[typed]{an} empty category
A  (defining \xout{An} the ``empty'' topos A\^{} )  as a precontractor,  although
formally (in terms of yesterday's definition) it is.  Thus I guess I'
ll introduce the following

\uline{Definition}   A topos is called \uline{vertical} if it is not an ``empty'' topos
(i.e. the category of sheaves on it is not equivalent to  the final
category \hand{$\Delta$}$_{\mathrm{o}}$ ),  and if moreover any \xout{subobject of the} open subtopos
is either the ``empty'' or the ``full'' one (hence the same for any sub-%
topos,  whether open or not).  A \xout{category} small category A is called
vertical,  if the associated topos (with category of sheaves A\^{}) is,

\margin{34}
